\section{Contribution to the Dissertation}
\label{sec:doctrine-contribution}

\subsection{Contribution to Prediction vs.~Coordination}
\label{sec:doctrine-contribution-prediction}

The central argument of the dissertation distinguishes prediction from coordination as the
two competing governance paradigms for consequential software. Prediction-based systems
observe, model, and forecast; coordination-based systems admit, refuse, receipt-bear, and
replay. The doctrine corpus is the formal statement of the coordination paradigm.

The 2016 language-model lesson that initiates the thesis lineage is precisely a demonstration
of the limits of prediction: local text imitation produces statistically plausible outputs
without conserving consequence. Prediction without admission gate produces narration, not
process truth. The Blue River Dam doctrine formalizes this gap and resolves it: ``a claim
that skips any step is narration. A claim that completes every step is process intelligence.''

The doctrine corpus contributes to the prediction vs.\ coordination argument by providing the
formal vocabulary for coordination: typed admission (\texttt{Admit::admit()}), named refusal
(\texttt{Refusal\textless{}R,W\textgreater{}}), typed witness (the \texttt{W} phantom marker),
and receipt-bearing execution (\texttt{Receipt\textless{}T,W\textgreater{}} with BLAKE3
chaining). These are not software patterns; they are formal coordination primitives. The
dissertation uses these primitives to argue that governance of consequential software requires
coordination-capable infrastructure, and that prediction-capable infrastructure --- however
sophisticated --- cannot substitute for it.

The MAPE-K autonomic actuation doctrine extends this contribution: it shows how coordination
can be made self-managing. The loop's Execute phase emits a receipt; the Knowledge base
accumulates only receipted actuations; the Plan phase is bounded to the elastic subnet
authorized by the compliance boundary. This is autonomous coordination under formal law, not
autonomous prediction under statistical inference.

\subsection{Contribution to Process Evidence}
\label{sec:doctrine-contribution-evidence}

The doctrine corpus makes the most direct contribution to the dissertation's process evidence
framework. It establishes the full evidential chain from raw event records to board-admissible
claims, and it specifies precisely what breaks the chain and how each breakage is detected.

The \texttt{Evidence\textless{}T,State,W\textgreater{}} carrier, the seven evidence lifecycle
states, the \texttt{Admit::admit()} gate, and the three failure modes (raw laundering, hidden
loss, unsigned projection) constitute a complete formal theory of process evidence. This
theory is grounded in published process mining literature: van der Aalst's OCEL 2.0 standard
(2023) provides the object-centric evidence model; Carmona et al.\ (2018) provide the
conformance checking framework; Berti and van der Aalst (2019) provide the token replay
algorithm; Adriansyah et al.\ (2011) provide the alignment-based conformance algorithm.

The doctrine corpus's contribution is to embed this literature in a type-theoretic framework
that enforces the evidential requirements at compile time. The fitness bound
\texttt{Between01\textless{}NUM,DEN\textgreater{}} enforced via \texttt{Require\textless{}\{NUM
<= DEN\}\textgreater{}: IsTrue} is not just a software invariant; it is the type-theoretic
translation of the formal requirement that fitness lie in $[0,1]$ as specified by Rozinat
and van der Aalst (2008). The \texttt{ConformanceViolation} with \texttt{StrictViolation::law()}
returning a \texttt{\&'static str} is the type-theoretic translation of the Van der Aalst
Constitution's dictum: ``model-vs-log mismatch is not a discrepancy; it is a defect.''

The process evidence contribution is also honest about its limits. The PARTIAL status
acknowledges that the evidence chain is specified in doctrine but not yet fully executed in
wasm4pm. The doctoral claim is not that the execution authority is operational; it is that
the formal specification of what operational execution requires is complete and paper-grounded.
This is a contribution to the theory of process evidence, with the execution as future
receipted work.

\subsection{Contribution to Manufacturing Architecture}
\label{sec:doctrine-contribution-manufacturing}

The doctrine corpus provides the foundational law for the CodeManufactory architecture that
spans this dissertation's project ecosystem. The manufacturing surface combinatorial formula
$\text{Surface} = |S| \times |L| \times |A| \times |F| \times |R| \times |Q| \times |B|
\times |D|$ defines the scope of the manufacturing problem: every combination of algorithm
family, log type, admissibility state, fitness metric, refusal reason, quality metric, board
claim type, and lifecycle stage is a distinct manufacturing surface that must be covered by
the type system.

The six algorithms defined in doctrine --- Paper-to-Law, PM4Py Oracle Mapping, Admission
Gate, Receipt-Bearing Execution, M\&A Deck Manufacturing, and Decommissioning --- constitute
the formal manufacturing pipeline. Each algorithm has a specified input type, output type,
gate condition, and receipt requirement. This specification is the authority layer from which
the \texttt{ggen} manufacturing machinery derives its projection targets and the wasm4pm
execution authority derives its algorithm intake requirements.

The \texttt{DOWNSTREAM\_AUTHORIZATION\_LAW.md} is the architectural mechanism by which the
doctrine corpus governs manufacturing: no downstream manufacturing run is authorized without
tracing to a sealed checkpoint. This prevents unauthorized manufacturing drift --- the
production of artifacts that cannot be receipted against doctrine, replayed against a known
log, or falsified by an adversarial auditor.

\subsection{Contribution to Capital-Flow Infrastructure}
\label{sec:doctrine-contribution-capital}

The thesis lineage concludes at capital-flow, settlement, and DAO infrastructure. The
doctrine corpus's contribution to this destination is the establishment of the refusal capital
concept and the board-admissibility criteria that make process evidence M\&A-ready.

Refusal capital is defined as $\Sigma \; \text{verified\_refusals\_with\_receipts}$ --- the
accumulated set of typed, witnessed, receipted refusals that the process intelligence system
has issued. This is described as the non-copyable competitive moat: an organization whose
process intelligence system has issued thousands of typed refusals with receipts has built an
audit history that a buyer, regulator, or counterparty can verify independently. This is
categorically different from an organization that has issued thousands of log entries.

The six board-admissibility criteria established in \texttt{MA\_READY\_PROCESS\_INTELLIGENCE.md}
translate directly into capital-flow infrastructure requirements: every board projection
$B_i$ must satisfy Evidence$_i$, TypeLaw$_i$, Receipt$_i$, Replay$_i$, Standard$_i$, and
Lifecycle$_i$. A settlement or DAO transaction that requires process evidence as a
precondition can use these criteria as the formal gate. This makes the doctrine corpus the
foundational specification layer for process-evidence-grounded capital allocation ---
the terminus of the thesis lineage from the 2016 language-model lesson to industry-complete
architecture.
